% =====================================================================
%  Lecture 06 -- E-UTRAN: the LTE Radio Access Network   (Persian edition)
% =====================================================================
\documentclass[11pt,a4paper]{article}
\providecommand\ACNroot{../..}
\input{\ACNroot/templates/lecture-fa}

\ACNset{number}{06}
\ACNset{title}{\lr{E-UTRAN}: شبکهٔ دسترسی رادیویی \lr{LTE}}
\ACNset{shorttitle}{\lr{E-UTRAN}}
\ACNset{titleother}{E-UTRAN: the LTE Radio Access Network}
\ACNset{part}{بخش دوم -- دسترسی و هستهٔ نسل چهارم (درس‌های ۵ تا ۱۳)}
\ACNset{time}{۷ تا ۹ ساعت (۳٫۵ ساعت مطالعه، ۴ تا ۵ ساعت تمرین)}
\ACNset{level}{متوسط}
\ACNset{prereq}{درس ۵؛ لایه‌بندی \lr{OSI}؛ مبانی \lr{ARQ}}
\ACNset{gen}{4}

\begin{document}
\ACNcover

% ---------------------------------------------------------------------
\section{اهداف یادگیری}
پس از مطالعهٔ این درس باید بتوانید:
\begin{objectives}
  \item کارکردهای \acro{eNB} را \textbf{برشمارید} و توضیح دهید چرا \acro{E-UTRAN} تخت توانست
        \acro{RNC} را در خود جذب کند.
  \item پشته‌های پروتکل صفحهٔ کنترل و صفحهٔ کاربر \acro{LTE} را \textbf{رسم کنید} و نقش
        \lr{PHY, MAC, RLC, PDCP} و \acro{RRC} را بیان کنید.
  \item کانال‌های منطقی را در هر دو جهت به کانال‌های انتقال و فیزیکی \textbf{نگاشت کنید}.
  \item واسط‌های \lr{S1} و \acro{X2}، پشته‌های پروتکل و رویه‌های اصلی آن‌ها را \textbf{توضیح دهید}.
  \item ماشین حالت \acro{RRC}، \acro{DRX} و افزونهٔ \lr{RRC\_INACTIVE} در \acro{NR} را \textbf{توصیف کنید}
        و مصالحهٔ توان/تأخیر \acro{DRX} را \textbf{کمّی کنید}.
  \item شماره‌های ترتیب \acro{PDCP} و پنجره‌های \acro{RLC} را برای یک نرخ دادهٔ هدف \textbf{اندازه‌گذاری کنید}.
\end{objectives}

% ---------------------------------------------------------------------
\section{دانش پیش‌نیاز}
این درس روی نیمهٔ چپ معماری \acro{EPS} درس ۵ تمرکز می‌کند. به ایدهٔ لایه‌بندی پروتکل (هر لایه به
لایهٔ بالاتر خدمت می‌دهد و یک سرآیند می‌افزاید)، مفهوم درخواست تکرار خودکار (\acro{ARQ}) با شماره‌های
ترتیب و پنجره‌ها، و مفهوم حامل \acro{EPS} نیاز دارید.

\paragraph{واژه‌های کلیدی} واژگان این درس همراه با معادل انگلیسی آن‌ها:
\begin{keyterms}
شبکهٔ دسترسی رادیویی \acro{LTE} & \lr{E-UTRAN} & شبکهٔ دسترسی \acro{LTE}: شبکه‌ای از \acro{eNB}ها بدون کنترل‌گر. \\
کنترل منابع رادیویی & \lr{RRC} & کنترل لایهٔ ۳ میان \acro{UE} و \acro{eNB}: اتصال، اندازه‌گیری، امنیت. \\
پروتکل همگرایی دادهٔ بسته‌ای & \lr{PDCP} & فشرده‌سازی سرآیند، رمزنگاری، یکپارچگی، تحویل به‌ترتیب. \\
کنترل پیوند رادیویی & \lr{RLC} & قطعه‌بندی و \acro{ARQ} در حالت‌های \lr{TM}، \lr{UM} یا \lr{AM}. \\
کنترل دسترسی به رسانه & \lr{MAC} & تسهیم، زمان‌بندی، \acro{HARQ}، دسترسی تصادفی. \\
کانال منطقی & \lr{Logical channel} & \emph{چه چیزی} حمل می‌شود (کنترل یا ترافیک). \\
کانال انتقال & \lr{Transport channel} & \emph{چگونه} میان \acro{MAC} و \acro{PHY} حمل می‌شود (قالب، کدگذاری). \\
کانال فیزیکی & \lr{Physical channel} & منابع واقعی زمان--فرکانس. \\
پروتکل‌های کاربردی \lr{S1} و \lr{X2} & \lr{S1AP / X2AP} & سیگنالینگ \lr{eNB--MME} و \lr{eNB--eNB} روی \acro{SCTP}. \\
دریافت ناپیوسته & \lr{DRX} & چرخه‌های خواب که باتری را به بهای تأخیر ذخیره می‌کنند. \\
\end{keyterms}

% ---------------------------------------------------------------------
\section{مقدمه و انگیزه}
در \acro{UMTS} پروتکل‌های رادیویی میان دو گره تقسیم شده بودند: \lr{NodeB} لایهٔ فیزیکی را و \acro{RNC}
همهٔ لایه‌های بالاتر را مدیریت می‌کرد. بنابراین هر تصمیم زمان‌بندی، هر ارسال مجدد \acro{RLC} و هر تحویل از
شبکهٔ انتقال \acro{Iub} عبور می‌کرد که تأخیر می‌افزود و \acro{RNC} را به گلوگاه تبدیل می‌کرد. \acro{HSPA}
زمان‌بندی سریع و \acro{HARQ} را به \lr{NodeB} برده بود. \acro{LTE} این حرکت را کامل کرد: \acro{eNB} \emph{همهٔ}
پروتکل‌های رادیویی را پایان می‌دهد؛ پس فاصلهٔ میان یک رویداد رادیویی و یک تصمیم درون یک جعبه می‌ماند.

این درس آن جعبه را باز می‌کند. دادهٔ کاربر و سیگنالینگ را در پشتهٔ پروتکل دنبال می‌کنیم، می‌بینیم انواع
اطلاعات چگونه در کانال‌ها سازمان می‌یابند و دو واسط شبکه‌ای \acro{eNB} را بررسی می‌کنیم: \lr{S1} به سوی
هسته و \acro{X2} به سوی همسایه‌ها. با ماشین حالت \acro{RRC} پایان می‌دهیم که تعیین می‌کند \acro{UE} و شبکه
چه مقدار باید در حافظه نگه دارند و یک گوشی بیکار با چه سرعتی می‌تواند ارسال داده را آغاز کند.

% ---------------------------------------------------------------------
\section{مفاهیم بنیادی}

\subsection{چرا \lr{RNC} نه؟}
\acro{RNC} به سه دلیل وجود داشت: ترکیب در تحویل نرم، مدیریت متمرکز منابع رادیویی بسیاری از \lr{NodeB}ها،
و پایان دادن به پروتکل‌هایی که برای ایستگاه‌های پایهٔ دههٔ ۱۹۹۰ بیش از حد پیچیده بودند. \acro{LTE} هر
دلیل را حذف کرد: \acro{OFDMA} از تحویل سخت استفاده می‌کند و ترکیب لازم نیست؛ زمان‌بندی هر ۱ میلی‌ثانیه
باید کنار رادیو انجام شود؛ و توان پردازش ارزان شده بود. هماهنگی میان سلول‌ها که پیش‌تر \acro{RNC} فراهم
می‌کرد با سیگنالینگ مستقیم \acro{eNB} به \acro{eNB} روی \acro{X2} و با کارکردهای شبکهٔ خودسازمان‌ده جایگزین شد.

\subsection{لایه‌بندی و خدمت حامل}
هر لایهٔ پشتهٔ رادیویی خدمتی کاملاً تعریف‌شده به بالا ارائه می‌دهد:
\begin{itemize}
  \item \textbf{\acro{PHY}} \emph{کانال‌های انتقال} را به \acro{MAC} ارائه می‌دهد.
  \item \textbf{\acro{MAC}} \emph{کانال‌های منطقی} را به \acro{RLC} ارائه می‌دهد.
  \item \textbf{\acro{RLC}} \emph{کانال‌های \lr{RLC}} را به \acro{PDCP} ارائه می‌دهد.
  \item \textbf{\acro{PDCP}} \emph{حامل‌های رادیویی} را به \acro{IP} (حامل‌های رادیویی داده، \acro{DRB}) و به
        \acro{RRC} (حامل‌های رادیویی سیگنالینگ، \acro{SRB}) ارائه می‌دهد.
\end{itemize}
یک حامل رادیویی داده بخش رادیویی یک حامل \acro{EPS} است (درس ۵)؛ \lr{SRB0} پیام‌های \acro{RRC} را پیش از
فعال شدن امنیت و \lr{SRB1} و \lr{SRB2} پیام‌های \acro{RRC} و \acro{NAS} را پس از آن حمل می‌کنند.

% ---------------------------------------------------------------------
\section{توضیحات فنی تفصیلی}

\subsection{کارکردهای \lr{eNB}}
\acro{eNB} این کارکردها را انجام می‌دهد \cite{ts36300}:
\begin{itemize}
  \item \textbf{مدیریت منابع رادیویی}: کنترل حامل رادیویی، کنترل پذیرش رادیویی، کنترل تحرک اتصال
        (تحویل)، تخصیص پویای منابع به \acro{UE}ها در فراسو و فروسو (زمان‌بندی)، هماهنگی تداخل میان سلولی
        (\acro{ICIC}) و توازن بار.
  \item \textbf{فشرده‌سازی سرآیند و رمزنگاری} دادهٔ کاربر (\acro{PDCP}).
  \item \textbf{انتخاب \acro{MME}} هنگام اتصال (تابع انتخاب گرهٔ \acro{NAS}) و مسیریابی دادهٔ کاربر به سوی \lr{S-GW}.
  \item \textbf{فراخوانی} و \textbf{پخش اطلاعات سامانه} که به ترتیب از \acro{MME} و بخش بهره‌برداری و
        نگهداری سرچشمه می‌گیرند.
  \item \textbf{پیکربندی اندازه‌گیری} برای تحرک و زمان‌بندی.
\end{itemize}

\subsection{لایه‌های پروتکل}
\paragraph{لایهٔ فیزیکی (\lr{PHY})} \acro{OFDMA} در فروسو و \acro{SC-FDMA} در فراسو با زیرحامل‌های
\qty{15}{\kilo\hertz}؛ کدگذاری توربو برای داده؛ \lr{QPSK}، \lr{16QAM} و \lr{64QAM} (\lr{256QAM} در نسخه‌های
بعدی)؛ تا \lr{MIMO} $4\times4$ (بعدها $8\times8$)؛ ترکیب نرم \acro{HARQ}؛ جست‌وجوی سلول با سیگنال‌های
همگام‌سازی اولیه و ثانویه (\lr{PSS/SSS}). شبکهٔ منابع و عدد‌شناسی آن موضوع درس ۷ است.

\paragraph{\lr{MAC}} کانال‌های منطقی را به کانال‌های انتقال نگاشت می‌کند، چند کانال منطقی را در یک بلوک
انتقال تسهیم می‌کند، تصمیم \emph{زمان‌بندی} را در هر \acro{TTI} می‌گیرد، \acro{HARQ} را اجرا می‌کند (۸ فرایند
موازی توقف-و-انتظار در \acro{FDD})، اولویت کانال‌های منطقی، دسترسی تصادفی (درس ۸)، پیش‌فرستی زمانی و
گزارش‌های وضعیت بافر و حاشیهٔ توان را مدیریت می‌کند \cite{ts36321}.

\paragraph{\lr{RLC}} سه حالت دارد \cite{ts36322}:
\begin{itemize}
  \item \textbf{حالت شفاف (\lr{TM})}: بدون سرآیند و بدون قطعه‌بندی -- برای پخش، فراخوانی و \lr{SRB0}.
  \item \textbf{حالت بدون تأیید (\lr{UM})}: قطعه‌بندی/الحاق و بازچینی، بدون ارسال مجدد -- برای ترافیک
        حساس به تأخیر مانند \acro{VoLTE}.
  \item \textbf{حالت با تأیید (\lr{AM})}: \acro{ARQ} با گزارش‌های وضعیت را می‌افزاید -- برای ترافیک \acro{TCP} و
        سیگنالینگ، جایی که قابلیت اطمینان از تأخیر مهم‌تر است.
\end{itemize}
\acro{RLC} بسته‌های \acro{IP} را با اندازهٔ بلوک انتقالی که زمان‌بند انتخاب کرده تطبیق می‌دهد: ممکن است یک
بسته را به قطعه‌هایی ببُرد یا چند بسته را در یک \acro{PDU} جای دهد.

\paragraph{\lr{PDCP}} \cite{ts36323} فشرده‌سازی مقاوم سرآیند (\acro{ROHC}) سرآیندهای \lr{IP/UDP/RTP}،
رمزنگاری دادهٔ کاربر و کنترل، حفاظت از یکپارچگی دادهٔ کنترل (\acro{SRB}) و تحویل به‌ترتیب و بدون تکرار
در طول تحویل (با انتقال وضعیت \acro{SN} درس ۴) را انجام می‌دهد. هر \acro{PDU} در \acro{PDCP} یک شمارهٔ ترتیب
(\acro{SN}) دارد؛ این شماره همراه با شمارهٔ ابرقاب (\acro{HFN}) که در هر دو سو نگه داشته می‌شود \emph{\lr{COUNT}}
سی‌ودو بیتی را می‌سازد، ورودی الگوریتم رمزنگاری که هرگز نباید زیر یک کلید تکرار شود.

\paragraph{\lr{RRC}} \cite{ts36331} اطلاعات سامانه را پخش می‌کند (\acro{MIB} روی \acro{PBCH}، \acro{SIB}ها روی
\acro{PDSCH})، اتصال \acro{RRC} را مدیریت می‌کند (برقراری، پیکربندی مجدد، برقراری مجدد، آزادسازی)،
اندازه‌گیری‌ها را پیکربندی و گزارش‌ها را دریافت می‌کند، امنیت لایهٔ دسترسی را فعال می‌کند و پیام‌های
\acro{NAS} را منتقل می‌کند.

\subsection{کانال‌ها}
\Cref{fig:channels} نگاشت اصلی را نشان می‌دهد. کانال‌های منطقی عبارت‌اند از کانال پخش (\acro{BCCH})،
فراخوانی (\acro{PCCH})، کنترل مشترک (\acro{CCCH}، پیش از وجود اتصال \acro{RRC})، کنترل اختصاصی (\acro{DCCH})
و ترافیک اختصاصی (\acro{DTCH}). کانال‌های انتقال آن‌ها را بر اساس قالب ارسال گروه‌بندی می‌کنند: کانال‌های
مشترک فروسو و فراسو (\lr{DL-SCH, UL-SCH}) تقریباً همه‌چیز را حمل می‌کنند. کانال‌های فیزیکی شامل کانال‌های
دادهٔ مشترک \acro{PDSCH} و \acro{PUSCH}، کانال‌های کنترلی بدون کانال انتقال (\acro{PDCCH} برای تخصیص‌های
زمان‌بندی، \acro{PCFICH} برای اندازهٔ ناحیهٔ کنترل، \acro{PHICH} برای تأییدهای \acro{HARQ} فراسو، \acro{PUCCH}
برای کنترل فراسو) و سیگنال‌های مرجع (\lr{CRS, DMRS, SRS}) برای اندازه‌گیری و آشکارسازی هستند.

\begin{acnfigure}
  \resizebox{\linewidth}{!}{\acnfig{lecture-06/channel-mapping}}
  \caption{نگاشت اصلی کانال‌های \acro{LTE}. تقریباً همهٔ کانال‌های منطقی به کانال‌های مشترک می‌رسند؛
  \lr{PDCCH, PHICH, PCFICH} و \acro{PUCCH} کنترل لایهٔ ۱/۲ را بدون کانال انتقال حمل می‌کنند.}
  \label{fig:channels}
\end{acnfigure}

\subsection{واسط \lr{S1}}
\lr{S1} هر \acro{eNB} را به یک یا چند \acro{MME} (\lr{S1-MME}) و \lr{S-GW} (\lr{S1-U}) وصل می‌کند -- رابطه‌ای
چندبه‌چند به نام \emph{\lr{S1-flex}} که تجمیع \acro{MME} را ممکن می‌کند (درس ۱۲).
\begin{itemize}
  \item \textbf{\lr{S1-MME}}: \lr{S1AP} روی \acro{SCTP} روی \acro{IP} \cite{ts36413}. \acro{SCTP} انتقالی قابل
        اطمینان و پیام‌محور با چندخانگی و جریان‌های متعدد فراهم می‌کند و از انسداد سرصف میان \acro{UE}ها
        جلوگیری می‌کند.
  \item \textbf{\lr{S1-U}}: \lr{GTP-U} روی \acro{UDP} روی \acro{IP}؛ یک تونل برای هر \lr{E-RAB} (بخش \lr{S1} یک حامل \acro{EPS}).
\end{itemize}
رویه‌های اصلی \lr{S1AP}: \emph{\lr{S1 Setup}} (\acro{eNB} و \acro{MME} هویت‌ها و نواحی ردیابی تحت پوشش را مبادله
می‌کنند)، \emph{\lr{Initial UE Message}} (نخستین پیام \acro{NAS} یک \acro{UE})، \emph{\lr{Initial Context Setup}}
(\acro{MME} زمینهٔ \acro{UE} را نصب می‌کند: کلید امنیتی، حامل‌ها، قابلیت‌ها)، \emph{\lr{E-RAB Setup/Modify/Release}}،
\emph{\lr{Paging}}، \emph{\lr{Handover Preparation/Resource Allocation}} (تحویل \lr{S1})، \emph{\lr{Path Switch}}،
\emph{\lr{UE Context Release}} و \emph{\lr{Downlink/Uplink NAS Transport}}.

\subsection{واسط \lr{X2}}
\acro{X2} \acro{eNB}های همسایه را به هم وصل می‌کند \cite{ts36423}: \lr{X2AP} روی \acro{SCTP} برای کنترل و
\lr{GTP-U} برای ارسال دادهٔ کاربر در حین تحویل. رویه‌های اصلی: \emph{\lr{X2 Setup}} (مبادلهٔ سلول‌های تحت
پوشش که روابط خودکار همسایگی را ممکن می‌کند)، \emph{\lr{Handover Preparation}}، \emph{\lr{SN Status Transfer}}،
\emph{\lr{UE Context Release}}، \emph{\lr{Load Indication}} (اطلاعات تداخل برای \acro{ICIC}) و
\emph{\lr{Resource Status Reporting}} (بار برای توازن بار تحرک). \acro{X2} اختیاری است: بدون آن، تحویل‌ها از
طریق \lr{S1} و \acro{MME} انجام می‌شوند.

\subsection{حالت‌های \lr{RRC} و \lr{DRX}}
\acro{LTE} دو حالت \acro{RRC} دارد (\cref{fig:rrc}):
\begin{itemize}
  \item \textbf{\lr{RRC\_IDLE}}: بدون اتصال \acro{RRC} و بدون زمینه در \acro{eNB}؛ \acro{UE} انتخاب مجدد سلول
        انجام می‌دهد، اطلاعات سامانه را می‌خواند و در هر چرخهٔ فراخوانی یک بار فراخوانی را پایش می‌کند
        (\acro{DRX} بیکار)؛ \acro{MME} آن را در سطح ناحیهٔ ردیابی می‌شناسد (\lr{ECM-IDLE}).
  \item \textbf{\lr{RRC\_CONNECTED}}: \acro{eNB} زمینهٔ \acro{UE} را نگه می‌دارد؛ \acro{UE} می‌تواند زمان‌بندی شود،
        اندازه‌گیری‌ها را گزارش می‌دهد و شبکه آن را تحویل می‌دهد؛ برای صرفه‌جویی در توان می‌تواند از
        \acro{DRX} حالت متصل استفاده کند و در هر چرخهٔ \acro{DRX} برای یک \emph{دورهٔ روشن} بیدار شود.
\end{itemize}
رفتن از \lr{IDLE} به \lr{CONNECTED} به دسترسی تصادفی، برقراری اتصال \acro{RRC}، درخواست خدمت \acro{NAS}،
فعال‌سازی امنیت و برقراری حامل نیاز دارد -- ده‌ها میلی‌ثانیه و دوازده پیام. بنابراین \acro{NR}
حالت \textbf{\lr{RRC\_INACTIVE}} را می‌افزاید \cite{ts38331}: \acro{UE} و \acro{gNB} زمینه (امنیت، حامل‌ها) را نگه
می‌دارند، هسته هنوز \acro{UE} را متصل می‌بیند و ازسرگیری تنها به یک مبادلهٔ کوتاه نیاز دارد.

\begin{acnfigure}
  \acnfig{lecture-06/rrc-states}
  \caption{ماشین حالت \acro{RRC}. \acro{LTE} حالت‌های \lr{IDLE} و \lr{CONNECTED} را دارد؛ \acro{NR} حالت
  \lr{INACTIVE} را می‌افزاید که در آن زمینه به‌جای آزاد شدن معلق می‌شود.}
  \label{fig:rrc}
\end{acnfigure}

% ---------------------------------------------------------------------
\section{مدل‌های ریاضی}

\subsection{بهرهٔ فشرده‌سازی سرآیند}\label{sec:rohc}
\paragraph{مسئله} \acro{ROHC} چه مقدار ظرفیت رادیویی برای \acro{VoLTE} صرفه‌جویی می‌کند؟
\paragraph{مدل} یک قاب گفتار $L_\mathrm{p}$ بایتی با $H$ بایت سرآیند \lr{RTP/UDP/IP} حمل می‌شود؛ \acro{ROHC}
در حالت پایدار $H$ را به سرآیند فشرده‌ای به اندازهٔ $H_\mathrm{c}$ بایت کاهش می‌دهد. بازده رادیویی بسته:
\begin{equation}\label{eq:rohc}
  \eta = \frac{L_\mathrm{p}}{L_\mathrm{p} + H}, \qquad
  G = \frac{L_\mathrm{p}+H}{L_\mathrm{p}+H_\mathrm{c}} \ \text{(capacity gain)}.
\end{equation}
\paragraph{مثال} $L_\mathrm{p} = 32$ بایت، $H = 40$ بایت (\lr{IPv4} ۲۰ + \acro{UDP} ۸ + \acro{RTP} ۱۲) و
$H_\mathrm{c} = 3$ بایت: $\eta$ از $44\%$ به $91\%$ می‌رسد و $G = 72/35 = 2.06$. با \lr{IPv6} ($H = 60$)
بهره $92/35 = 2.6$ است. \acro{ROHC} تعداد کاربران \acro{VoLTE} قابل حمل در یک سلول را تقریباً دو برابر می‌کند؛
روی \lr{S1-U} اعمال نمی‌شود و سرآیندها در آنجا کامل حمل می‌شوند (درس ۵).
\paragraph{محدودیت‌ها} اندازهٔ فشردهٔ حالت پایدار تنها پس از برقراری یک زمینه به دست می‌آید؛ اتلاف بسته
می‌تواند فشرده‌گشا را ناهمگام کند؛ بنابراین نمایه‌های \acro{ROHC} فشرده‌سازی را با استحکام معاوضه می‌کنند.

\subsection{فضای شمارهٔ ترتیب و چرخش آن}\label{sec:sn}
\paragraph{مسئله} در یک نرخ معین چند بیت \acro{SN} برای \acro{PDCP} لازم است؟
\paragraph{مدل} یک پروتکل پنجره‌ای با \acro{SN} $n$ بیتی حداکثر می‌تواند $2^{n-1}$ \acro{PDU} معوق داشته باشد
(نصف فضا)، وگرنه \acro{PDU}های قدیمی و جدید مبهم می‌شوند. با نرخ $\lambda_\mathrm{p}$ \acro{PDU} در ثانیه،
پنجره این زمان را پوشش می‌دهد:
\begin{equation}\label{eq:snwin}
  T_{\mathrm{win}} = \frac{2^{\,n-1}}{\lambda_\mathrm{p}}, \qquad
  \lambda_\mathrm{p} = \frac{R}{8L},
\end{equation}
برای نرخ بیت $R$ و \acro{PDU}های $L$ بایتی. $T_{\mathrm{win}}$ باید از بیشترین تأخیر بازچینی یا ارسال مجدد،
که معمولاً شامل \acro{HARQ}، \acro{ARQ} در \acro{RLC} و ارسال دادهٔ تحویل است، بیشتر باشد.
\paragraph{مثال} $R = \qty{300}{\mega\bit\per\second}$ و $L = 1500$ بایت: $\lambda_\mathrm{p} = 25\,000$
\acro{PDU} بر ثانیه. \acro{SN} دوازده‌بیتی $T_{\mathrm{win}} = 2048/25\,000 = 82$ میلی‌ثانیه می‌دهد -- در حین
تحویل به‌طور ناخوشایندی کوتاه. به همین دلیل \acro{LTE} \acro{SN} پانزده‌بیتی (۶۵۵ میلی‌ثانیه) و \acro{NR} \acro{SN}
هجده‌بیتی \acro{PDCP} را افزودند.

\subsection{گذردهی حالت با تأیید \lr{RLC}}\label{sec:rlcwin}
برای پنجره‌ای از $W$ \acro{PDU} با $L$ بایت و زمان رفت‌وبرگشت وضعیت $T_{\mathrm{RTT}}$، فرستنده‌ای که محدود
به پنجره است نمی‌تواند از این مقدار فراتر رود:
\begin{equation}\label{eq:rlc}
  R_{\max} = \frac{8\,W\,L}{T_{\mathrm{RTT}}} .
\end{equation}
با \acro{SN} ده‌بیتی \lr{RLC AM} در \acro{LTE}، $W = 512$؛ برای $L = 1500$ بایت و $T_{\mathrm{RTT}} = 20$ میلی‌ثانیه،
$R_{\max} = 512\times12\,000/0.02 = \qty{307}{\mega\bit\per\second}$: برای نسخهٔ ۸ کافی است و دلیل \acro{SN}های
طولانی‌تر در نسخه‌های بعدی با نرخ‌های بالاتر همین است (\cref{exr:rlc}).

\subsection{\lr{DRX}: توان در برابر تأخیر}\label{sec:drx}
\paragraph{مسئله} یک \acro{UE} متصل چه مدت باید بخوابد؟
\paragraph{مدل} چرخهٔ \acro{DRX} برابر $T_{\mathrm{DRX}}$ با دورهٔ روشن $T_{\mathrm{on}}$؛ توان $P_{\mathrm{on}}$ در بیداری
و $P_{\mathrm{sleep}}$ در خواب؛ بسته‌های فروسو در لحظه‌های تصادفی می‌رسند (در غیر این صورت در دورهٔ خواب
ترافیکی نیست).
\paragraph{استخراج} با چرخهٔ کار $d = T_{\mathrm{on}}/T_{\mathrm{DRX}}$،
\begin{equation}\label{eq:drx}
  \bar P = d\,P_{\mathrm{on}} + (1-d)\,P_{\mathrm{sleep}}, \qquad
  \bar D_{\mathrm{DRX}} \approx \frac{T_{\mathrm{DRX}}}{2},
\end{equation}
زیرا بسته‌ای که در زمانی تصادفی و یکنواخت می‌رسد به‌طور میانگین نصف چرخه را برای دورهٔ روشن بعدی صبر
می‌کند.
\paragraph{مثال و تفسیر} $T_{\mathrm{on}} = 10$ میلی‌ثانیه، $P_{\mathrm{on}} = 500$ میلی‌وات و
$P_{\mathrm{sleep}} = 10$ میلی‌وات (نمایشی): چرخهٔ ۴۰ میلی‌ثانیه‌ای ۱۳۲٫۵ میلی‌وات و ۲۰ میلی‌ثانیه تأخیر
اضافی می‌دهد؛ چرخهٔ ۳۲۰ میلی‌ثانیه‌ای ۲۵٫۳ میلی‌وات و ۱۶۰ میلی‌ثانیه؛ ۲۵۶۰ میلی‌ثانیه ۱۱٫۹ میلی‌وات اما ۱٫۲۸
ثانیه (\cref{fig:drx}). توان مانند $1/T_{\mathrm{DRX}}$ کاهش می‌یابد تا $P_{\mathrm{sleep}}$ غالب شود، در حالی که
تأخیر خطی رشد می‌کند: فراتر از چندصد میلی‌ثانیه، چرخه‌های طولانی‌تر انرژی اندکی را به بهای تأخیر زیاد
می‌خرند. دستگاه‌های واقعی پس از هر بسته به مدت یک زمان‌سنج بی‌فعالیتی بیدار می‌مانند که مدل نادیده می‌گیرد.

\begin{acnfigure}
  \acnfig{lecture-06/drx-power}
  \caption{توان میانگین (محور چپ) و میانگین تأخیر افزوده (محور راست) بر حسب چرخهٔ \acro{DRX} برای
  $T_{\mathrm{on}} = 10$\,\lr{ms}، بر اساس \cref{eq:drx}.}
  \label{fig:drx}
\end{acnfigure}

% ---------------------------------------------------------------------
\section{معماری شبکه و نمودارها}
\Cref{fig:stacks} پشته‌های رادیویی را خلاصه می‌کند. در صفحهٔ کنترل، \acro{RRC} در \acro{eNB} پایان می‌یابد در
حالی که \acro{NAS} به‌طور شفاف از آن عبور می‌کند و به \acro{MME} می‌رسد (به سوی \acro{UE} درون \acro{RRC} و به
سوی \acro{MME} درون \lr{S1AP}). در صفحهٔ کاربر، \acro{PDCP}، \acro{RLC} و \acro{MAC} در \acro{eNB} پایان می‌یابند
و \acro{eNB} بسته‌های \acro{IP} را به تونل‌های \lr{GTP-U} روی \lr{S1-U} رله می‌کند.

\begin{acnfigure}
  \resizebox{\linewidth}{!}{\acnfig{lecture-06/protocol-stacks}}
  \caption{پشته‌های پروتکل رادیویی \acro{LTE}: صفحهٔ کنترل (چپ) و صفحهٔ کاربر (راست).}
  \label{fig:stacks}
\end{acnfigure}

% ---------------------------------------------------------------------
\section{رویه‌های گام‌به‌گام}
\begin{procedurebox}{از \lr{RRC\_IDLE} تا انتقال داده (آغازشده از موبایل)}\label{proc:idle2conn}
\begin{steps}
  \item \textbf{دسترسی تصادفی.} \acro{UE} یک پیش‌نوا می‌فرستد و پاسخ دسترسی تصادفی را همراه با پیش‌فرستی
        زمانی و مجوز دریافت می‌کند (درس ۸).
  \item \textbf{درخواست اتصال \acro{RRC}.} روی \lr{SRB0} پیام \msg{RRC Connection Request} (هویت و علت) را
        می‌فرستد؛ \acro{eNB} با \msg{RRC Connection Setup} پاسخ می‌دهد و \lr{SRB1} ساخته می‌شود.
  \item \textbf{انتقال \acro{NAS}.} \acro{UE} پیام \msg{RRC Connection Setup Complete} را همراه با \acro{NAS}
        \msg{Service Request} می‌فرستد؛ \acro{eNB} آن را در \lr{S1AP} \msg{Initial UE Message} به \acro{MME}ای که
        از روی هویت \acro{UE} انتخاب کرده ارسال می‌کند.
  \item \textbf{برقراری زمینه.} \acro{MME} با \lr{S1AP} \msg{Initial Context Setup Request} همراه با کلید امنیتی
        $K_{\mathrm{eNB}}$، \lr{E-RAB}ها و نقاط پایانی تونل \lr{S-GW} پاسخ می‌دهد.
  \item \textbf{امنیت لایهٔ دسترسی.} \acro{eNB} پیام \msg{Security Mode Command} را می‌فرستد؛ \acro{UE} کلیدهای
        لایهٔ دسترسی را استخراج و با \msg{Security Mode Complete} پاسخ می‌دهد. اکنون \acro{PDCP} رمزنگاری و
        حفاظت از یکپارچگی را انجام می‌دهد.
  \item \textbf{پیکربندی حامل.} \acro{eNB} پیام \msg{RRC Connection Reconfiguration} را می‌فرستد که \lr{SRB2} و
        \acro{DRB}ها را می‌سازد؛ پس از \msg{Reconfiguration Complete} با \msg{Initial Context Setup Response} همراه
        با نقاط پایانی تونل \lr{S1-U} خودش پاسخ می‌دهد. دادهٔ فراسو و فروسو جریان می‌یابد.
\end{steps}
\end{procedurebox}

% ---------------------------------------------------------------------
\section{مثال‌های حل‌شده}

\begin{examplebox}{سربار پروتکل یک قطعهٔ \lr{TCP}}\label{ex:overhead}
\textbf{مسئله.} یک محمولهٔ \acro{TCP} ۱۴۶۰ بایتی روی \acro{LTE} با \lr{IPv4} ارسال می‌شود. بایت‌های سرآیندی را
که روی رادیو توسط \acro{PDCP} (۲ بایت، \acro{SN} دوازده‌بیتی)، \lr{RLC AM} (۲ بایت، بدون قطعه‌بندی) و \acro{MAC}
(زیرسرآیند ۲ بایتی) و روی \lr{S1-U} افزوده می‌شود بشمارید.

\textbf{حل.} بستهٔ \acro{IP}: $1460 + 20 + 20 = 1500$ بایت. رادیو: $1500 + 2 + 2 + 2 = 1506$ بایت (۰٫۴٪ سربار؛
\acro{ROHC} معمولاً برای \acro{TCP} به کار نمی‌رود). \lr{S1-U}: $1500 + 36 = 1536$ بایت (۲٫۴٪).

\textbf{تفسیر.} برای بسته‌های بزرگ سربار لایهٔ ۲ ناچیز است؛ هزینه‌های واقعی سیگنال‌های مرجع و کانال‌های
کنترلی \acro{PHY} (درس ۵) و ارسال‌های مجددند.
\end{examplebox}

\begin{examplebox}{پنجرهٔ \lr{PDCP} در حین تحویل}\label{ex:pdcp}
\textbf{مسئله.} یک \acro{UE} با نرخ \lr{150\,Mbit/s} و بسته‌های ۱۵۰۰ بایتی داده دریافت می‌کند. در حین تحویل،
بسته‌های ارسال‌شده ممکن است تا ۶۰ میلی‌ثانیه دیر برسند. آیا \acro{SN} دوازده‌بیتی \acro{PDCP} کافی است؟

\textbf{حل.} $\lambda_\mathrm{p} = 150\times10^6/12\,000 = 12\,500$ \acro{PDU} بر ثانیه. طبق \cref{eq:snwin}:
$T_{\mathrm{win}} = 2048/12\,500 = 164$ میلی‌ثانیه $> 60$. بله، با حاشیه؛ در \lr{300\,Mbit/s} (۸۲ میلی‌ثانیه)
حاشیه اندک می‌شد.
\end{examplebox}

\begin{examplebox}{انتخاب چرخهٔ \lr{DRX}}\label{ex:drx}
\textbf{مسئله.} یک برنامهٔ پیام‌رسان به‌طور میانگین ۲۰۰ میلی‌ثانیه تأخیر اضافی را تحمل می‌کند. با پارامترهای
\cref{sec:drx}، طولانی‌ترین چرخهٔ استاندارد \acro{DRX} از میان ۴۰، ۸۰، ۱۶۰ و ۳۲۰ میلی‌ثانیه را انتخاب و
صرفه‌جویی توان را نسبت به حالت بدون \acro{DRX} (پیوسته $P_{\mathrm{on}}$) محاسبه کنید.

\textbf{حل.} $T_{\mathrm{DRX}}/2 \le 200$ میلی‌ثانیه، ۳۲۰ میلی‌ثانیه را مجاز می‌کند (۱۶۰ میلی‌ثانیه تأخیر).
$\bar P = 25.3$ میلی‌وات در برابر ۵۰۰ میلی‌وات: کاهش ۹۵ درصدی.

\textbf{تفسیر.} بیشتر صرفه‌جویی از قبل در ۱۶۰ میلی‌ثانیه (۴۰٫۶ میلی‌وات) به دست می‌آید؛ اپراتور می‌تواند برای
هر برنامه از طریق زمان‌سنج بی‌فعالیتی و تنظیمات چرخه انتخاب کند.
\end{examplebox}

% ---------------------------------------------------------------------
\section{مقایسهٔ فناوری‌ها}
\begin{table}[htbp]
\centering\small
\caption{کارکردهای شبکهٔ دسترسی رادیویی: \acro{UTRAN} در برابر \acro{E-UTRAN} در برابر \acro{NG-RAN}.}\label{tab:ran}
\rowcolors{2}{acnmist}{white}
\begin{tabularx}{\linewidth}{@{}l Y Y Y@{}}
\toprule
\rowcolor{white} & \acro{UTRAN} (نسل سوم) & \acro{E-UTRAN} (نسل چهارم) & \acro{NG-RAN} (نسل پنجم) \\
\midrule
گره‌ها              & \lr{NodeB + RNC}            & \acro{eNB}                   & \acro{gNB} (در صورت نیاز \lr{CU + DU}) \\
محل \lr{RLC/PDCP}   & \acro{RNC}                  & \acro{eNB}                   & \acro{DU} (\acro{RLC})، \acro{CU} (\acro{PDCP}) \\
زمان‌بند            & \acro{RNC} (\lr{R99})، \lr{NodeB} (\acro{HSPA}) & \acro{eNB}، ۱ میلی‌ثانیه & \lr{gNB-DU}، شیار/ریزشیار \\
پیوند میان گره‌ها   & \lr{Iur}                    & \lr{X2}                      & \lr{Xn}؛ \lr{F1} میان \lr{CU} و \lr{DU} \\
واسط هسته           & \lr{Iu (RANAP)}             & \lr{S1 (S1AP)}               & \lr{NG (NGAP)} \\
حالت‌های \lr{RRC}   & \lr{Idle, Cell\_DCH/FACH/PCH, URA\_PCH} & \lr{Idle, Connected} & \lr{Idle, Inactive, Connected} \\
زیرلایهٔ جدید لایهٔ ۲ & --                        & --                           & \acro{SDAP} (نگاشت جریان کیفیت خدمات) \\
\bottomrule
\end{tabularx}
\end{table}

\paragraph{مصالحه} قرار دادن همه‌چیز در \acro{eNB} تأخیر را کمینه کرد اما هماهنگی میان سلول‌ها (تداخل، ارسال
مشترک) را دشوارتر کرد، زیرا اکنون به سیگنالینگ \acro{X2} میان گره‌های هم‌تراز نیاز دارد. \lr{5G} با تقسیم
\lr{CU/DU} شبکهٔ دسترسی را تا حدی دوباره متمرکز می‌کند (درس ۱۴): لایه‌های پایینی حساس به زمان در سایت
آنتن می‌مانند و \acro{PDCP} و \acro{RRC} را می‌توان تجمیع کرد -- سازشی آگاهانه میان تمرکز \acro{UTRAN} و تخت
بودن \acro{E-UTRAN}.

% ---------------------------------------------------------------------
\section{نکته‌های کلیدی}
\begin{takeawaybox}
\begin{itemize}
  \item \acro{eNB} لایه‌های \lr{PHY, MAC, RLC, PDCP} و \acro{RRC} را پایان می‌دهد: تصمیم‌های رادیویی در یک گره می‌مانند.
  \item \acro{PDCP}: \acro{ROHC}، رمزنگاری، یکپارچگی (\acro{SRB}ها)، تحویل به‌ترتیب؛ \lr{COUNT} $=$ \lr{HFN} $\Vert$ \lr{SN} هرگز نباید تکرار شود.
  \item \acro{RLC}: \lr{TM}، \lr{UM} (صدا)، \lr{AM} (\acro{TCP}، سیگنالینگ)؛ \acro{MAC}: تسهیم، زمان‌بندی، \acro{HARQ}.
  \item کانال‌های منطقی $\to$ انتقال $\to$ فیزیکی؛ تقریباً همه‌چیز از کانال‌های مشترک استفاده می‌کند.
  \item \lr{S1-MME} (\lr{S1AP/SCTP}) و \lr{S1-U} (\lr{GTP-U/UDP}) به سوی هسته؛ \acro{X2} میان \acro{eNB}ها.
  \item \lr{RRC\_IDLE} در برابر \lr{RRC\_CONNECTED} (+ \lr{RRC\_INACTIVE} در \acro{NR})؛ \acro{DRX} توان
        ($\propto1/T_{\mathrm{DRX}}$) را با تأخیر ($T_{\mathrm{DRX}}/2$) معاوضه می‌کند.
  \item فضای \acro{SN} باید زمان بازچینی را بپوشاند: $2^{n-1}/\lambda_\mathrm{p} > T_{\mathrm{reorder}}$.
\end{itemize}
\end{takeawaybox}

% ---------------------------------------------------------------------
\section{تمرین‌ها}

\begin{exercisebox}{1}{کدام لایه؟}\label{exr:layer}
برای هر کارکرد، لایه‌ای از \acro{LTE} را که آن را انجام می‌دهد نام ببرید: \acro{HARQ}، \acro{ARQ}، رمزنگاری دادهٔ
کاربر، قطعه‌بندی، فشرده‌سازی سرآیند، حفاظت از یکپارچگی پیام‌های \acro{RRC}، انتخاب \acro{MCS}، پیکربندی اندازه‌گیری.
\end{exercisebox}

\begin{exercisebox}{1}{چرا \lr{SCTP}؟}\label{exr:sctp}
توضیح دهید چرا \lr{S1AP} و \lr{X2AP} به‌جای \acro{TCP} یا \acro{UDP} از \acro{SCTP} استفاده می‌کنند. یک سناریوی
خرابی را مثال بزنید که چندخانگی آن را مدیریت می‌کند.
\end{exercisebox}

\begin{exercisebox}{2}{پنجرهٔ \lr{RLC}}\label{exr:rlc}
با \cref{eq:rlc}، پنجرهٔ \lr{RLC AM} لازم برای \lr{1\,Gbit/s} با \acro{PDU}های ۱۵۰۰ بایتی و زمان رفت‌وبرگشت
وضعیت ۱۵ میلی‌ثانیه و حداقل تعداد بیت‌های \acro{SN} را بیابید.
\end{exercisebox}

\begin{exercisebox}{2}{\lr{ROHC} و \lr{IPv6}}\label{exr:rohc}
\cref{sec:rohc} را برای \lr{AMR-WB} با نرخ \lr{23.85\,kbit/s} (قاب‌های ۶۱ بایتی در هر ۲۰ میلی‌ثانیه) و
\lr{IPv6} تکرار کنید. نرخ بیت روی رادیو با و بدون \acro{ROHC} و بهرهٔ ظرفیت چقدر است؟
\end{exercisebox}

\begin{exercisebox}{3}{\lr{DRX} با زمان‌سنج بی‌فعالیتی}\label{exr:drxinact}
پس از هر بستهٔ دریافتی، \acro{UE} به مدت زمان‌سنج بی‌فعالیتی $T_{\mathrm{in}}$ بیدار می‌ماند. بسته‌ها به‌صورت
فرایند پواسون با نرخ $\lambda$ می‌رسند. با فرض $T_{\mathrm{in}} \ll 1/\lambda$ عبارتی تقریبی برای کسر زمان
بیداری استخراج و برای $\lambda = 0.2$\,\lr{s}$^{-1}$، $T_{\mathrm{in}} = 100$ میلی‌ثانیه، $T_{\mathrm{on}} = 10$
میلی‌ثانیه و $T_{\mathrm{DRX}} = 320$ میلی‌ثانیه ارزیابی کنید.
\end{exercisebox}

\begin{exercisebox}{3}{بیکار در برابر متصل}\label{exr:idleconn}
آزاد کردن \acro{UE} به \lr{RRC\_IDLE} پس از بی‌فعالیتی، حافظهٔ \acro{eNB} و توان \acro{UE} را ذخیره می‌کند، اما هر
انفجار دادهٔ جدید حدود ۱۲ پیام سیگنالینگ هزینه دارد. اگر \acro{eNB} بتواند $10^4$ گذار از این نوع در ثانیه
را مدیریت کند و هر \acro{UE} هر ۳۰ ثانیه یک انفجار تولید کند، از نظر سیگنالینگ صرف به چند \acro{UE} می‌توان
خدمت داد؟ بحث کنید \lr{RRC\_INACTIVE} پاسخ را چگونه تغییر می‌دهد.
\end{exercisebox}

\subsection*{پاسخ‌ها}
\addcontentsline{toc}{subsection}{پاسخ‌ها}

\begin{solutionbox}{\ref{exr:layer}}
\acro{HARQ}: \acro{MAC}. \acro{ARQ}: \acro{RLC} (\lr{AM}). رمزنگاری دادهٔ کاربر: \acro{PDCP}. قطعه‌بندی: \acro{RLC}.
فشرده‌سازی سرآیند: \acro{PDCP}. یکپارچگی \acro{RRC}: \acro{PDCP}. انتخاب \acro{MCS}: زمان‌بند \acro{MAC} (با ورودی
تطبیق پیوند از \acro{PHY}). پیکربندی اندازه‌گیری: \acro{RRC}.
\end{solutionbox}

\begin{solutionbox}{\ref{exr:sctp}}
\acro{SCTP} پیام‌محور است (پیام‌های سیگنالینگ مرزهای خود را حفظ می‌کنند)، مانند \acro{TCP} قابل اطمینان و دارای
کنترل ازدحام است، از جریان‌های متعدد پشتیبانی می‌کند تا پیام از دست رفتهٔ یک \acro{UE} پیام‌های دیگران را
مسدود نکند (بدون انسداد سرصف میان جریان‌ها) و از چندخانگی پشتیبانی می‌کند: یک ارتباط می‌تواند از چند مسیر
\acro{IP} استفاده کند. اگر یک پیوند انتقال میان \acro{eNB} و \acro{MME} قطع شود، ارتباط از نشانی دیگر ادامه
می‌یابد بی‌آنکه واسط \lr{S1} و همراه با آن همهٔ زمینه‌های \acro{UE} از بین بروند.
\end{solutionbox}

\begin{solutionbox}{\ref{exr:rlc}}
$W \ge R\,T_{\mathrm{RTT}}/(8L) = 10^9\times0.015/12\,000 = 1250$ \acro{PDU}. فضای \acro{SN} باید دست‌کم
$2W = 2500$ باشد؛ پس ۱۲ بیت ($4096$) لازم است؛ ده بیت نسخهٔ ۸ \acro{LTE} کافی نیست (از این رو شماره‌های ترتیب
طولانی‌تر \acro{RLC} در نسخه‌های بعدی؛ \lr{RLC AM} در \acro{NR} از ۱۲ یا ۱۸ بیت استفاده می‌کند).
\end{solutionbox}

\begin{solutionbox}{\ref{exr:rohc}}
\lr{IPv6/UDP/RTP}: $40+8+12 = 60$ بایت. بدون \acro{ROHC}: $61+60 = 121$ بایت در هر ۲۰ میلی‌ثانیه $=$
\lr{48.4\,kbit/s}؛ با \acro{ROHC} (۳ بایت): ۶۴ بایت $=$ \lr{25.6\,kbit/s}. بهره $121/64 = 1.89$. بهره از قاب‌های ۳۲
بایتی کمتر است، زیرا محموله بزرگ‌تر است.
\end{solutionbox}

\begin{solutionbox}{\ref{exr:drxinact}}
هر بسته \acro{UE} را حدود $T_{\mathrm{in}}$ بیدار نگه می‌دارد؛ با $\lambda T_{\mathrm{in}} \ll 1$ این دوره‌ها به‌ندرت
هم‌پوشانی دارند. در بقیهٔ زمان \acro{DRX} اعمال می‌شود. کسر بیداری
$f \approx \lambda T_{\mathrm{in}} + (1-\lambda T_{\mathrm{in}})\,T_{\mathrm{on}}/T_{\mathrm{DRX}}
= 0.02 + 0.98\times0.03125 = 0.0506$، یعنی ۵٫۱٪. با توان‌های \cref{sec:drx}،
$\bar P \approx 0.0506\times500 + 0.9494\times10 = 34.8$ میلی‌وات. زمان‌سنج بی‌فعالیتی برای این ترافیک حدود ۴۰٪
به انرژی حالت \acro{DRX} خالص می‌افزاید.
\end{solutionbox}

\begin{solutionbox}{\ref{exr:idleconn}}
هر \acro{UE} در هر ثانیه $1/30$ گذار ایجاد می‌کند؛ پس $N = 10^4\times30 = 3\times10^5$ \acro{UE} برای هر
\acro{eNB} -- بیش از ظرفیت هر سلول، اما \acro{MME} مجموع بسیاری از \acro{eNB}ها را می‌بیند و پیام‌های فراخوانی،
\acro{NAS} و \lr{S1} هسته را نیز بار می‌کنند. با \lr{RRC\_INACTIVE} ازسرگیری تنها حدود ۳ پیام رادیویی و هیچ
سیگنالینگ هسته‌ای لازم دارد (هسته \acro{UE} را متصل می‌بیند)؛ بنابراین بار هسته حذف و بار رادیویی حدود ۴ برابر
کم می‌شود، به بهای ذخیرهٔ زمینه‌ها در شبکهٔ دسترسی.
\end{solutionbox}

% ---------------------------------------------------------------------
\section{مطالعهٔ بیشتر}
توصیف کلی \acro{E-UTRAN} \cite{ts36300} بهترین نقطهٔ شروع است؛ مشخصات لایه‌ها \acro{MAC} \cite{ts36321}، \acro{RLC}
\cite{ts36322}، \acro{PDCP} \cite{ts36323} و \acro{RRC} \cite{ts36331} هستند؛ واسط‌های شبکه \lr{S1AP} \cite{ts36413}
و \lr{X2AP} \cite{ts36423} و \acro{SCTP} در \cite{rfc9260} مشخص شده است. \lr{Sesia} و همکاران
\cite{sesia2011lte} چند فصل را به پروتکل‌های لایهٔ ۲ اختصاص داده‌اند.

\begin{latin}
\printbibliography[heading=subbibliography,title={References}]
\end{latin}
\end{document}
